\section{Viable sets for the ground state of a gapped chain}
\label{sec:lvv_viable_sets}

Landau, Vazirani and Vidick build an approximation to the ground state
$\Gamma$ of a gapped chain of $L$ sites cut by cut, keeping at the cut after
$i$ sites a small set $S$ of vectors on the first $i$ sites that supports the
left half of a good approximation to
$\Gamma$~\cite[lines 240--244]{Landau2015Polynomial}. For a vector $\psi$ on
the chain, write $M_i(\psi)$ for its coefficient matrix across the cut after
$i$ sites, with rows indexed by the configurations of the first $i$ sites.
The reduced state of $\psi$ on the first $i$ sites is $M_i(\psi)M_i(\psi)^*$,
whose range is the span of the columns of $M_i(\psi)$.

\begin{definition}[Viable sets]
    \label{def:lvv_viable_set}
    \lean{TNLean.MPS.AreaLaw.cutColumn,
        TNLean.MPS.AreaLaw.leftSupportSubmodule,
        TNLean.MPS.AreaLaw.mem_leftSupportSubmodule,
        TNLean.MPS.AreaLaw.IsViable}
    \leanok
    \uses{def:chain_model}
    Let $\Gamma$ be a vector on the chain, $\delta\in\mathbb R$ and $S$ a set of
    vectors on the first $i$ sites. Then $S$ is $(i,\delta)$-viable for $\Gamma$
    if some unit vector $\phi$ with $|\langle\phi,\Gamma\rangle|\ge1-\delta$ has
    every column of $M_i(\phi)$ in $\operatorname{Span}(S)$, that is, $\phi$ lies
    in $\operatorname{Span}(S)\otimes\mathcal H_R$ and its reduced state on the
    first $i$ sites is supported on $\operatorname{Span}(S)$. This is
    Definition~def:viable of~\cite[lines 240--242]{Landau2015Polynomial}.
\end{definition}

\begin{theorem}[Support through the projection onto the span]
    \label{thm:lvv_viable_projection}
    \lean{TNLean.MPS.AreaLaw.spanProjection,
        TNLean.MPS.AreaLaw.isStarProjection_spanProjection,
        TNLean.MPS.AreaLaw.isHermitian_spanProjection,
        TNLean.MPS.AreaLaw.spanProjection_mul_self,
        TNLean.MPS.AreaLaw.spanProjection_mul_cutMatrix_eq_iff,
        TNLean.MPS.AreaLaw.regionState_leftBlock_eq,
        TNLean.MPS.AreaLaw.spanProjection_mul_regionState_eq_iff,
        TNLean.MPS.AreaLaw.leftApply_cutVector_eq_iff,
        TNLean.MPS.AreaLaw.leftApply_spanProjection_cutVector_eq_iff}
    \leanok
    \uses{def:lvv_viable_set}
    Let $P_S$ be the orthogonal projection onto $\operatorname{Span}(S)$, a
    Hermitian idempotent matrix on the configurations of the first $i$ sites.
    For every vector $\psi$ the following are equivalent: the columns of
    $M_i(\psi)$ lie in $\operatorname{Span}(S)$; $P_SM_i(\psi)=M_i(\psi)$;
    $P_S\rho=\rho$ for the reduced state $\rho=M_i(\psi)M_i(\psi)^*$ of $\psi$ on
    the first $i$ sites; $(P_S\otimes I)\psi=\psi$, with $\psi$ written in the
    coordinates of the cut after the first $i$ sites.
\end{theorem}

\begin{proof}
    \leanok
    A column $c$ of $M_i(\psi)$ lies in $\operatorname{Span}(S)$ exactly when
    $P_Sc=c$. The range of $M_i(\psi)M_i(\psi)^*$ is the column space of
    $M_i(\psi)$. In the coordinates of the cut, $\psi$ is $M_i(\psi)$ read as a
    vector, and $(P\otimes I)\psi$ is $PM_i(\psi)$, for every matrix $P$.
\end{proof}

\begin{theorem}[Viable sets and viable subspaces]
    \label{thm:lvv_viable_closeness}
    \lean{TNLean.MPS.AreaLaw.IsViable.le_norm_starProjection,
        TNLean.MPS.AreaLaw.isViable_iff_le_norm_starProjection,
        TNLean.MPS.AreaLaw.IsViable.forall_mem_span_singleton,
        TNLean.MPS.AreaLaw.isViable_iff_forall_mem_span_singleton}
    \leanok
    \uses{def:lvv_viable_set}
    Let $U=\operatorname{Span}(S)\otimes\mathcal H_R$ and let $P_U$ be the
    orthogonal projection onto $U$. If $S$ is $(i,\delta)$-viable for $\Gamma$,
    then $\|P_U\Gamma\|\ge1-\delta$, and for $\delta<1$ the converse holds. For a
    unit vector $\Gamma$ and $\delta\le1$, viability gives
    $(1-\delta')\|t\|^2\le\|P_Ut\|^2$ for every $t$ in the span of $\Gamma$, with
    $\delta'=1-(1-\delta)^2$, and for $\delta<1$ the converse holds. This is the
    closeness of $U$ to the span of $\Gamma$ in the sense of the viable subspaces
    of Arad, Landau, Vazirani and
    Vidick~\cite[lines 433--440]{Arad2016RigorousRG}, with $\delta'$ in place of
    $\delta$, since $\delta$ bounds an overlap and $\delta'$ a squared overlap.
\end{theorem}

\begin{proof}
    \leanok
    For a unit $\phi\in U$, $|\langle\phi,\Gamma\rangle|
        =|\langle\phi,P_U\Gamma\rangle|\le\|P_U\Gamma\|$, with equality for
    $\phi=P_U\Gamma/\|P_U\Gamma\|$ when $P_U\Gamma\ne0$. Squaring gives the second
    form.
\end{proof}

\begin{definition}[Viable sets of matrix product states]
    \label{def:lvv_viable_mps}
    \lean{TNLean.MPS.AreaLaw.leftConfigOfFin,
        TNLean.MPS.AreaLaw.IsLeftOpenMPS,
        TNLean.MPS.AreaLaw.IsViableMPS}
    \leanok
    \uses{def:lvv_viable_set, def:chain_open_mps}
    The set $S$ is $(i,s,b,\delta)$-viable for $\Gamma$ if it is
    $(i,\delta)$-viable, has at most $s$ members, and each member is an open
    matrix product state on the first $i$ sites of bond dimension at most
    $\max(1,b)$~\cite[line 243]{Landau2015Polynomial}. Bond dimension $0$ is read
    as $1$; the source allows $b=0$ only for the initial set on no
    sites~\cite{gap:lvv13_bond_dimension_zero}.
\end{definition}

\begin{theorem}[Extending a viable set by one site]
    \label{thm:lvv_viable_extension}
    \lean{TNLean.MPS.AreaLaw.IsOpenMPS.snoc,
        TNLean.MPS.AreaLaw.restrictLeft,
        TNLean.MPS.AreaLaw.extendSite,
        TNLean.MPS.AreaLaw.lastLeftSite,
        TNLean.MPS.AreaLaw.extendSite_apply,
        TNLean.MPS.AreaLaw.extend,
        TNLean.MPS.AreaLaw.extendSite_mem_extend,
        TNLean.MPS.AreaLaw.card_extend_le,
        TNLean.MPS.AreaLaw.consRightConfig,
        TNLean.MPS.AreaLaw.cutColumn_eq_sum_extendSite,
        TNLean.MPS.AreaLaw.leftSupportSubmodule_le_extend,
        TNLean.MPS.AreaLaw.IsViable.extend,
        TNLean.MPS.AreaLaw.IsLeftOpenMPS.extendSite,
        TNLean.MPS.AreaLaw.IsViableMPS.extend}
    \leanok
    \uses{def:lvv_viable_mps}
    Let $1\le i\le L$, let $q$ be the local dimension, and let $S'$ be the set
    of vectors $v\otimes|k\rangle$ on the first $i$ sites, for $v\in S$ on the
    first $i-1$ sites and $k$ a basis state. If $S$ is
    $(i-1,s,b,\delta)$-viable for $\Gamma$, then $S'$ is
    $(i,qs,b,\delta)$-viable for $\Gamma$, with the same witness. This is
    Claim~c:ext of~\cite[lines 379--400]{Landau2015Polynomial}, stated there
    for $\delta=c_\varepsilon/n$.
\end{theorem}

\begin{proof}
    \leanok
    \uses{def:lvv_viable_mps}
    $S'$ has at most $q|S|$ members. Let $w$ be an open matrix product state on
    $n$ sites with bond dimension at most $b'=\max(1,b)$, and let $f$ be a
    vector on one site. For a cut after $j<n$ sites,
    \[
        M_j(w\otimes f)=M_j(w)\,R_j,\qquad
        (R_j)_{z',(z,x)}=\delta_{z'z}\,f(x),
    \]
    so $\operatorname{rank}M_j(w\otimes f)\le\operatorname{rank}M_j(w)\le b'$,
    while the new cut after $n$ sites gives $M_n(w\otimes f)=wf^{\mathsf T}$, of
    rank at most $1\le b'$. With $f=|k\rangle$, every member of $S'$ has bond
    dimension at most $b'$. For the witness, decomposing site $i$ on the basis
    gives, for each configuration $\beta$ of the sites after $i$,
    \[
        M_i(\phi)_{\cdot,\beta}=\sum_k M_{i-1}(\phi)_{\cdot,(k,\beta)}\otimes|k\rangle.
    \]
    When the columns of $M_{i-1}(\phi)$ lie in $\operatorname{Span}(S)$, each
    summand lies in $\operatorname{Span}(S')$, so the columns of $M_i(\phi)$ do as
    well; the witness $\phi$ is unchanged.
\end{proof}

\begin{theorem}[Bond dimensions in a finite span]
    \label{thm:open_mps_finite_span}
    \lean{TNLean.MPS.AreaLaw.isOpenMPS_of_mem_span}
    \leanok
    \uses{def:chain_open_mps}
    Let $L,q,D$ be natural numbers, and let $S$ be a finite set of vectors
    on a chain of $L$ sites of local dimension $q$. Suppose that every
    member of $S$ is an open-boundary matrix product state of bond
    dimension at most $D$. Every $v\in\operatorname{Span}(S)$ is then an
    open-boundary matrix product state of bond dimension at most $|S|D$.
    This is the finite-span estimate used
    in~\cite[Claim~2.7, source lines 520--527]{Landau2015Polynomial}.
\end{theorem}
\begin{proof}\leanok
    \uses{thm:chain_cut_rank_of_mps}
    Write $v=\sum_{w\in S}a_w w$. At every internal cut $j$,
    \begin{align}
        M_j(v)                    & =\sum_{w\in S}a_wM_j(w),\notag \\
        \operatorname{rank}M_j(v) & \le
        \sum_{w\in S}\operatorname{rank}(a_wM_j(w))\le |S|D.\notag
    \end{align}
    The cut-rank characterization gives the asserted representation,
    including empty and one-site chains and $D=0$.
\end{proof}

\begin{theorem}[Bond dimensions in a left span]
    \label{thm:left_mps_finite_span}
    \lean{TNLean.MPS.AreaLaw.isLeftOpenMPS_of_mem_span}
    \leanok
    \uses{def:lvv_viable_mps}
    Fix natural numbers $L,q,k,b$, and put $\beta=\max\{1,b\}$.
    Write $Rv$ for the vector read on $k$ sites by restriction to the left
    block, so $Rv(a)=v(a|_{\{j<L:j<k\}})$ for a $k$-site configuration $a$.
    Let $S$ be a finite set of left-block vectors, and suppose that every
    $Rw$, $w\in S$, is a matrix product state of bond dimension at most
    $\beta$. For every $v\in\operatorname{Span}(S)$, the vector $Rv$ is a
    matrix product state of bond dimension at most $\max\{1,|S|\beta\}$.
    When $S$ is nonempty and $b\ge1$, this bound is $|S|b$.
\end{theorem}
\begin{proof}\leanok
    \uses{thm:open_mps_finite_span,thm:chain_cut_rank_of_mps}
    The map $R$ is linear, so $Rv\in\operatorname{Span}(R(S))$ and
    $|R(S)|\le |S|$. Theorem~\ref{thm:open_mps_finite_span} gives a bond
    bound $|R(S)|\beta\le |S|\beta\le\max\{1,|S|\beta\}$.
\end{proof}

\begin{definition}[Left columns of a finite family]
    \label{def:lvv_finite_left_columns}
    \lean{TNLean.MPS.AreaLaw.leftColumns}
    \leanok
    \uses{def:chain_model}
    Let $I$ be a finite index set and let $m$ be a natural number. Write
    $\mathcal L$ for the Hilbert space of configurations on the left block
    of a cut. For vectors $u_i\in\mathcal L\otimes\mathbb C^m$, write
    $u_i=\sum_{j<m}u_{i,j}\otimes|j\rangle$ in the standard basis of the
    bond space, and define
    \begin{align}
        T=\{u_{i,j}:i\in I,\ j<m\}.\notag
    \end{align}
    This is the column construction used in the size-trimming step
    of~\cite[Claim~2.7, source lines 472--477 and 520--527]{Landau2015Polynomial}.
\end{definition}

\begin{theorem}[Number and support of the left columns]
    \label{thm:lvv_finite_left_columns}
    \lean{TNLean.MPS.AreaLaw.card_leftColumns_le,
        TNLean.MPS.AreaLaw.leftColumns_subset_span}
    \leanok
    \uses{def:lvv_finite_left_columns}
    The finite set $T$ satisfies $|T|\le |I|m$.
    Let $S\subseteq\mathcal L$, and let $P_S$ be the orthogonal projection
    onto $\operatorname{Span}(S)$. If $(P_S\otimes I_m)u_i=u_i$ for every
    $i\in I$, then $T\subseteq\operatorname{Span}(S)$.
\end{theorem}
\begin{proof}\leanok
    \uses{thm:lvv_viable_projection}
    The set $T$ is the image of $I\times\{0,\ldots,m-1\}$, so its cardinality
    is at most $|I|m$. Let $C_i$ have columns $u_{i,j}$. The fixed-vector
    identity gives $P_SC_i=C_i$, hence $P_Su_{i,j}=u_{i,j}$ for every
    retained column. The range of $P_S$ is $\operatorname{Span}(S)$.
\end{proof}


\begin{theorem}[Left support after extending the right factor]
    \label{thm:lvv_left_columns_right_extension}
    \lean{TNLean.MPS.AreaLaw.rightApply_mem_leftSupportSubmodule_leftColumns,
        TNLean.MPS.AreaLaw.rightApply_mem_leftSupportSubmodule_of_columns}
    \leanok
    \uses{def:lvv_finite_left_columns,def:lvv_viable_set}
    Let $T$ be the finite set of left columns in
    Definition~\ref{def:lvv_finite_left_columns}.
    For every $i\in I$ and every linear map $W:\mathbb C^m\to\mathcal H_R$,
    the chain vector obtained from $(I_{\mathcal L}\otimes W)u_i$ by the
    cut identification belongs to $\operatorname{Span}(T)\otimes\mathcal H_R$.
    More generally, an arbitrary finite environment and right extension
    preserve this support whenever every left column lies in
    $\operatorname{Span}(T)$.
    This is the support step used
    in~\cite[Claim~2.7, source lines 520--527]{Landau2015Polynomial}.
\end{theorem}
\begin{proof}\leanok
    Let $C_i$ be the coefficient matrix of $u_i$. The extended vector has
    coefficient matrix $C_iW^{\mathsf T}$. Its column at a right
    configuration $b$ is
    \begin{align}
        \sum_{j<m}W_{b,j}u_{i,j}.\notag
    \end{align}
    Each $u_{i,j}$ belongs to $T$, so every such column lies in
    $\operatorname{Span}(T)$.
\end{proof}

\begin{theorem}[A finite family of minimizing boundary states]
    \label{thm:lvv_finite_boundary_selection}
    \lean{TNLean.MPS.AreaLaw.exists_finite_boundary_minimizer_family}
    \leanok
    \uses{def:lvv_finite_left_columns,def:chain_gluing_left_coordinates,def:lvv_viable_set}
    Fix a set $S\subseteq\mathcal L$, a linear operator $H_L$ on $\mathcal L$,
    a natural number $m$, a real number $\eta$, and a finite set $N$ of
    matrices on $\mathcal C\otimes\mathbb C^m$, where $\mathcal C$ is the
    Hilbert space of the region between consecutive clipped left blocks. Identify
    $\mathcal L=\mathcal A\otimes\mathcal C$ by joining the left coordinates.
    Let $P_S$ be the orthogonal projection onto $\operatorname{Span}(S)$, and put
    \begin{align}
        Q       & =P_S\otimes I_m, \notag                                          \\
        E(\rho) & =\operatorname{Re}\tr\bigl(\rho(H_L\otimes I_m)\bigr),\notag     \\
        \mathcal F_X
                & =\left\{\rho\ge0:\tr\rho=1,\ Q\rho=\rho,\
        \bigl\|\operatorname{tr}_{\mathcal A}\rho-X\bigr\|_1\le\eta\right\},\notag \\
        F       & =\{X\in N:\mathcal F_X\ne\varnothing\}.\notag
    \end{align}
    For every $X\in F$, one can choose $\sigma_X\in\mathcal F_X$, a unit
    vector $u_X$ and a real number $\mu_X>0$ such that
    $E(\sigma_X)\le E(\rho)$ for every $\rho\in\mathcal F_X$,
    $\sigma_Xu_X=\mu_Xu_X$ and $Qu_X=u_X$.
    Every real eigenvalue of $\sigma_X$ with a nonzero eigenvector is at
    most $\mu_X$. After joining the left coordinates, write
    $u_X=\sum_{j<m}u_{X,j}\otimes|j\rangle$ and define
    \begin{align}
        T   & =\{u_{X,j}:X\in F,\ j<m\}.\notag       \\
        |T| & \le m|N|,\notag                        \\
        T   & \subseteq\operatorname{Span}(S).\notag
    \end{align}
    These choices use only $S,H_L,m,\eta,N$. This is an auxiliary
    fixed-bond construction for the size-trimming step
    of~\cite[Claim~2.7, source lines 472--478 and 520--527]{Landau2015Polynomial}.
\end{theorem}
\begin{proof}\leanok
    \uses{thm:chain_gluing_left_projection,thm:lvv_finite_left_columns}
    At every $X\in F$, choose a global minimizer and a leading unit
    eigenvector by the minimizing-density selection
    theorem~\cite[``A minimizing density matrix and a leading eigenvector'']{QICLean2026Blueprint}.
    Its positive eigenvalue gives $Qu_X=u_X$. Joining the left coordinates
    transports this identity to $(P_S\otimes I_m)u_X=u_X$.
    Theorem~\ref{thm:lvv_finite_left_columns} gives
    $T\subseteq\operatorname{Span}(S)$ and $|T|\le m|F|\le m|N|$.
\end{proof}


\begin{theorem}[A size-trimming family chosen before the ground state]
    \label{thm:lvv_size_trimmed_family}
    \lean{TNLean.MPS.AreaLaw.ChainHamiltonian.exists_size_trimmed_family,
        TNLean.MPS.AreaLaw.ChainHamiltonian.exists_size_trimmed_column_family}
    \leanok
    \uses{def:chain_model,def:chain_constant_bond,def:lvv_viable_set,def:lvv_viable_mps}
    Let $L\ge1$ and $q\ge2$. Fix a range-one chain Hamiltonian
    $H=\sum_Xh_X$ with positive semidefinite terms of norm at most one,
    a number $\Delta>0$, and a natural cut position $k$, with left blocks
    intersected with the chain. Let $S$ be a finite set of left-block
    vectors. Fix $b\in\mathbb N$, put $\beta=\max\{1,b\}$, and suppose that
    every member of $S$, read on $k$ sites as in
    Definition~\ref{def:lvv_viable_mps}, has matrix product state bond
    dimension at most $\beta$. Use Definition~\ref{def:chain_constant_bond}
    with $R=J=1$, and define
    \begin{align}
        c    & =(\Delta/169)^2,\notag                                   \\
        D    & =B_c,\notag                                              \\
        \eta & =c/(2L),\notag                                           \\
        K    & =\bigl(2\lceil(qD)^2/\eta\rceil+1\bigr)^{2(qD)^2}.\notag
    \end{align}
    There is a finite set $T$ of left-block vectors such that
    \begin{align}
        |T| & \le(D+1)DK,\notag                      \\
        T   & \subseteq\operatorname{Span}(S).\notag
    \end{align}
    Every member of $T$, read on $k$ sites, has matrix product state bond
    dimension at most $\max\{1,|S|\beta\}$. The set $T$ is chosen before
    both the ground-state energy and the ground vector: for every
    $E_0\in\mathbb R$ and every $\Gamma$, if $\Gamma$ is a normalized ground
    state of $H$ of energy $E_0$ with gap at least $\Delta$, then
    $(k,c/L)$-viability of $S$ for $\Gamma$ implies
    $(k,1/12)$-viability of $T$ for $\Gamma$.
    In this case there are $m\le D$, an isometry
    $W:\C^m\to\mathcal H_R$, and a unit vector
    $u\in\mathcal H_L\otimes\C^m$ such that every column
    $u_a=(I\otimes\langle a|)u$ belongs to $T$ and
    $|\langle\Gamma,(I\otimes W)u\rangle|\ge11/12$.

    This is an auxiliary version of
    \cite[Claim~2.7, source lines 472--478 and 520--527]{Landau2015Polynomial}.
    It uses the library bond constant, a union over actual bond dimensions,
    and the corrected matrix-space exponent. These qualifications are
    recorded in~\cite{gap:lvv13_constant_bond,gap:lvv13_trace_norm_net};
    the printed quantitative AKLV dependence remains separate.
\end{theorem}
\begin{proof}\leanok
    \uses{thm:chain_gluing_hamiltonian_decomposition,thm:chain_boundary_trace_norm_net,
        thm:lvv_finite_boundary_selection,thm:lvv_uniform_leading_bounds_library_bond,
        thm:chain_gluing_left_extension,thm:lvv_left_columns_right_extension,
        thm:left_mps_finite_span}
    First fix a decomposition $H=H_L+B$ with the stated left and right
    supports and $\|B\|\le2L$. For every integer $0\le m\le D$, choose an
    $\eta$-net $N_m$ on the actual boundary space, with $|N_m|\le K$.
    Theorem~\ref{thm:lvv_finite_boundary_selection} chooses a global
    minimizer and a leading unit eigenvector at exactly the feasible
    points of $N_m$. Let $T_m$ be their left columns, and put
    \begin{align}
        T   & =\bigcup_{m=0}^{D}T_m,\notag \\
        |T| & \le\sum_{m=0}^{D}|T_m|
        \le\sum_{m=0}^{D}m|N_m|\le(D+1)DK.\notag
    \end{align}
    All these choices precede $E_0$ and $\Gamma$. Every $T_m$ lies in
    $\operatorname{Span}(S)$, and Theorem~\ref{thm:left_mps_finite_span}
    gives the asserted bond bound.

    Now fix a ground-state pair $E_0,\Gamma$ for which $S$ is
    $(k,c/L)$-viable. Theorem~\ref{thm:lvv_uniform_leading_bounds_library_bond}
    supplies an actual bond $m\le D$, an isometry $W$ and a unit left
    state $x$ fixed by the left support projection. Write $\mathcal A$
    for the Hilbert space on $\lbrack0,\max\{k-1,0\})$, intersected with the chain. For
    $\rho_x=|x\rangle\langle x|$ and its boundary state $Y$, one has
    \begin{align}
        \rho_x                 & \ge0,\notag                                      \\
        \tr\rho_x              & =\|x\|^2=1,\notag                                \\
        (P_S\otimes I_m)\rho_x & =\rho_x,\notag                                   \\
        Y                      & =\operatorname{tr}_{\mathcal A}\rho_x\ge0,\notag \\
        \|Y\|_1                & =\tr Y=1.\notag
    \end{align}
    Choose $X\in N_m$ with $\|Y-X\|_1\le\eta$. The density $\rho_x$ is
    feasible at $X$, so $X$ belongs to the selected feasible set.
    The same uniform theorem applies to the already chosen global
    minimizer and its already chosen top eigenvector $u_{m,X}$. It gives
    a unit physical extension $v=(I\otimes W)u_{m,X}$ with
    $|\langle\Gamma,v\rangle|\ge11/12$.
    Joining the left coordinates and extending the right factor places
    $v$ in $\operatorname{Span}(T_m)\otimes\mathcal H_R$, hence in
    $\operatorname{Span}(T)\otimes\mathcal H_R$. Thus $v$ witnesses the
    asserted viability. The same selected eigenvector $u_{m,X}$, after
    joining the left coordinates, has all its columns in $T_m\subseteq T$;
    its norm is one and the original right extension $W$ is an isometry.
\end{proof}

\begin{theorem}[Polynomial size of the size-trimming family]
    \label{thm:lvv_size_trimmed_family_polynomial}
    \lean{TNLean.MPS.AreaLaw.ChainHamiltonian.exists_size_trimmed_family_polynomial,
        TNLean.MPS.AreaLaw.ChainHamiltonian.exists_size_trimmed_column_family_polynomial}
    \leanok
    \uses{def:chain_model,def:chain_constant_bond,def:lvv_viable_set,def:lvv_viable_mps}
    Under the assumptions of Theorem~\ref{thm:lvv_size_trimmed_family}, put
    \begin{align}
        c & =(\Delta/169)^2,\notag        \\
        D & =B_c,\notag                   \\
        p & =2(qD)^2,\notag               \\
        A & =\lceil2(qD)^2/c\rceil,\notag \\
        C & =(D+1)D(2A+1)^p.\notag
    \end{align}
    There is a finite set $T$, chosen before $E_0$ and $\Gamma$, such that
    \begin{align}
        |T| & \le CL^p,\notag                        \\
        T   & \subseteq\operatorname{Span}(S).\notag
    \end{align}
    Every member of $T$, read on $k$ sites, has matrix product state bond
    dimension at most $\max\{1,|S|\max\{1,b\}\}$.
    For every $E_0\in\mathbb R$ and every normalized ground state $\Gamma$
    of $H$ of energy $E_0$ with gap at least $\Delta$, $(k,c/L)$-viability of $S$
    implies $(k,1/12)$-viability of $T$.
    For fixed $q,\Delta$, both $C$ and $p$ are independent of $L$.
    This is the polynomial form of the auxiliary preceding theorem,
    with the same qualifications on the bond constant and net exponent.
\end{theorem}
\begin{proof}\leanok
    \uses{thm:lvv_size_trimmed_family,thm:chain_boundary_net_polynomial_bound}
    Take the same set $T$ as in Theorem~\ref{thm:lvv_size_trimmed_family}.
    Theorem~\ref{thm:chain_boundary_net_polynomial_bound} gives
    \begin{align}
        K   & \le(2A+1)^pL^p,\notag                      \\
        |T| & \le(D+1)DK\le(D+1)D(2A+1)^pL^p=CL^p.\notag
    \end{align}
    Its span inclusion, bond bound and viability implication are unchanged.
\end{proof}


\begin{theorem}[Counting words without singleton letters]
    \label{thm:lvv_even_word_count}
    \lean{Finset.card_image_le_half_of_no_singleton,
        Fintype.exists_factor_fin_of_card_range_le,
        Fintype.card_functions_of_card_range_le,
        Fintype.card_even_words_without_singleton}
    \leanok
    Let $p\ge1$ be an integer, and let $B$ be an alphabet with
    $\ell\ge1$ letters. The number of words of length $2p$ over $B$ in
    which no letter occurs exactly once is at most
    \begin{align}
        \ell^p p^{2p}.
        \label{eq:lvv_even_word_count}
    \end{align}
    This elementary count is one ingredient of an even-moment
    argument for sampled operators. It does not assert a matrix
    concentration inequality.
\end{theorem}
\begin{proof}\leanok
    Each occurring letter appears at least twice, so such a word
    uses at most $p$ letters. Every word with at most $p$ letters
    factors as a map from its $2p$ positions to a set of $p$ places,
    followed by a map from those places to $B$. Extend the second map
    arbitrarily on unused places. There are $p^{2p}\ell^p$ pairs of
    maps, which gives the bound.
\end{proof}


\begin{theorem}[Cancellation of a singleton sample label]
    \label{thm:lvv_singleton_word_cancellation}
    \lean{Fintype.sum_ordered_word_eq_zero_of_singleton}
    \leanok
    Let $A$ and $B$ be finite sets, and let $R$ be a possibly
    noncommutative semiring. Let $s:\{0,\ldots,m-1\}\to B$ specify
    the sample label at each position. Suppose that $s(t)$ occurs
    only at position $t$. Let $\Phi_k:A\to R$ satisfy
    $\sum_{a\in A}\Phi_t(a)=0$. Then
    \begin{align}
        \sum_{\sigma:B\to A}
        \Phi_0(\sigma(s(0)))\cdots\Phi_{m-1}(\sigma(s(m-1)))=0.
        \label{eq:lvv_singleton_word_cancellation}
    \end{align}
    The product is taken in the displayed order.
\end{theorem}
\begin{proof}\leanok
    Fix the values of $\sigma$ away from $s(t)$. The factors before
    and after position $t$ are then fixed, so the sum over
    $a=\sigma(s(t))$ has the form
    \begin{align}
        L\left(\sum_{a\in A}\Phi_t(a)\right)R'=0.
        \label{eq:lvv_singleton_word_cancellation_inner}
    \end{align}
    Sum this identity over the remaining values of $\sigma$.
\end{proof}


\begin{theorem}[Finite uniform moments of centered matrix samples]
    \label{thm:lvv_finite_uniform_matrix_moment}
    \lean{Matrix.average_norm_average_pow_le_of_centered}
    \leanok
    \uses{thm:lvv_even_word_count,thm:lvv_singleton_word_cancellation}
    Let $A,B$ be finite nonempty sets, with $|B|=\ell$, and let $N\ge1$.
    Let $X_b:A\to M_N(\mathbb C)$ satisfy
    \begin{align}
        \sum_{a\in A}X_b(a) & =0,   \\
        \|X_b(a)\|          & \le R
        \qquad(a\in A,\ b\in B),
        \label{eq:lvv_finite_moment_hypotheses}
    \end{align}
    where $R\ge0$ and the norm is the operator norm. For every integer $p\ge1$,
    \begin{align}
        \frac1{|A|^\ell}\sum_{\sigma:B\to A}
        \left\|\frac1\ell\sum_{b\in B}X_b(\sigma(b))\right\|^{2p}
        \le N\left(\frac{p^2R^2}{\ell}\right)^p.
        \label{eq:lvv_finite_uniform_matrix_moment}
    \end{align}
    The samples need not be Hermitian. This is an auxiliary finite
    moment estimate. The sharp matrix Bernstein inequality is proved
    separately in Theorem~\ref{thm:lvv_matrix_bernstein}.
\end{theorem}
\begin{proof}\leanok
    \uses{thm:lvv_even_word_count,thm:lvv_singleton_word_cancellation}
    Put $S_\sigma=\sum_bX_b(\sigma(b))$. Expand
    $\operatorname{tr}((S_\sigma^*S_\sigma)^p)$ into ordered words of length
    $2p$ in $B$, alternating adjointed and unadjointed factors.
    If a sample label occurs exactly once, its sum over $\sigma$ vanishes
    by Theorem~\ref{thm:lvv_singleton_word_cancellation}, since both
    $X_b$ and $X_b^*$ are centered. By Theorem~\ref{thm:lvv_even_word_count},
    at most $\ell^pp^{2p}$ words remain. Each has trace of absolute value
    at most $NR^{2p}$, because $|\operatorname{tr}T|\le N\|T\|$ and the
    operator norm is submultiplicative.

    The positive spectral theorem gives
    $\|S_\sigma\|^{2p}\le\operatorname{tr}((S_\sigma^*S_\sigma)^p)$.
    Sum this inequality over $\sigma$, apply the preceding bounds, and
    divide by $|A|^\ell\ell^{2p}$.
\end{proof}


\begin{theorem}[Finite uniform tails and centered sample existence]
    \label{thm:lvv_finite_uniform_matrix_tail}
    \lean{Matrix.finite_uniform_tail_bound_of_centered,
        Matrix.exists_average_norm_lt_of_centered,
        Matrix.exists_average_norm_lt_of_sample_card}
    \leanok
    \uses{thm:lvv_finite_uniform_matrix_moment}
    Under the hypotheses of Theorem~\ref{thm:lvv_finite_uniform_matrix_moment},
    let $\eta>0$. The fraction of assignments $\sigma:B\to A$ for which
    $\|\ell^{-1}\sum_bX_b(\sigma(b))\|\ge\eta$ is at most
    \begin{align}
        N\left(\frac{p^2R^2}{\ell\eta^2}\right)^p.
        \label{eq:lvv_finite_uniform_matrix_tail}
    \end{align}
    If this bound is less than one, an assignment with norm below $\eta$
    exists. In particular, it suffices that
    \begin{align}
        N          & <4^p,       \\
        \ell\eta^2 & \ge4p^2R^2.
        \label{eq:lvv_finite_uniform_centered_sample_count}
    \end{align}
\end{theorem}
\begin{proof}\leanok
    \uses{thm:lvv_finite_uniform_matrix_moment}
    Every assignment in the exceptional set contributes at least $\eta^{2p}$
    to the moment sum. Divide by the total number of assignments and
    apply Theorem~\ref{thm:lvv_finite_uniform_matrix_moment}. If the resulting
    fraction is less than one, the exceptional set is not the whole set.
    Under~\eqref{eq:lvv_finite_uniform_centered_sample_count}, the bound is
    at most $N4^{-p}<1$.
\end{proof}

\begin{theorem}[Approximating a finite matrix average by a sample]
    \label{thm:lvv_finite_matrix_sample_mean}
    \lean{Matrix.exists_average_sub_mean_norm_lt}
    \leanok
    \uses{thm:lvv_finite_uniform_matrix_tail}
    Let $A,B$ be finite nonempty sets, with $|B|=\ell$, and let $N\ge1$.
    Suppose that $Y:A\to M_N(\mathbb C)$ satisfies $\|Y(a)\|\le R$ for all
    $a\in A$, where $R\ge0$. Put $\mu=|A|^{-1}\sum_aY(a)$.
    If $\eta>0$ and an integer $p\ge1$ satisfies
    \begin{align}
        N          & <4^p,        \\
        \ell\eta^2 & \ge16p^2R^2,
        \label{eq:lvv_finite_uniform_sample_count}
    \end{align}
    there is an assignment $\sigma:B\to A$ for which
    \begin{align}
        \left\|\frac1\ell\sum_{b\in B}Y(\sigma(b))-\mu\right\|<\eta.
        \label{eq:lvv_finite_matrix_sample_mean}
    \end{align}
\end{theorem}
\begin{proof}\leanok
    \uses{thm:lvv_finite_uniform_matrix_tail}
    The triangle inequality gives $\|\mu\|\le R$. Thus $Y(a)-\mu$ is
    centered and has norm at most $2R$. Apply
    Theorem~\ref{thm:lvv_finite_uniform_matrix_tail} with $2R$ in place of $R$.
\end{proof}


\subsection{Exponential trace concavity}

\begin{theorem}[Unnormalized Klein inequality]
    \label{thm:matrix_unnormalized_klein}
    \lean{Matrix.trace_sub_le_quantumRelativeEntropy}
    \leanok
    For positive-definite complex matrices $X,Y$ of the same finite size,
    put $D_0(X\Vert Y)=\operatorname{Re}\operatorname{tr}(X(\log X-\log Y))$.
    Then
    \begin{align}
        \operatorname{Re}\operatorname{tr}X-
        \operatorname{Re}\operatorname{tr}Y\le D_0(X\Vert Y).
        \label{eq:matrix_unnormalized_klein}
    \end{align}
    This is~\cite[Fact~3]{Tropp2012JointConvexity}, with the affine trace
    correction separated from the relative entropy.
\end{theorem}
\begin{proof}\leanok
    In eigenbases of $X$ and $Y$, the squared overlaps form a doubly stochastic
    matrix. Sum the scalar inequality $x-y\le x(\log x-\log y)$ against these
    overlaps.
\end{proof}

\begin{theorem}[Attained variational formula for the trace]
    \label{thm:matrix_trace_variational_maximum}
    \lean{Matrix.isGreatest_trace_variational}
    \leanok
    \uses{thm:matrix_unnormalized_klein}
    For every positive-definite $Y$,
    \begin{align}
        \operatorname{Re}\operatorname{tr}Y
        =\max_{X>0}\operatorname{Re}\operatorname{tr}
        (X\log Y-X\log X+X).
        \label{eq:matrix_trace_variational_maximum}
    \end{align}
    The maximum is attained at $X=Y$. This is
    \cite[Lemma~6]{Tropp2012JointConvexity}.
\end{theorem}
\begin{proof}\leanok
    \uses{thm:matrix_unnormalized_klein}
    Rearrange Theorem~\ref{thm:matrix_unnormalized_klein} and take $X=Y$.
\end{proof}

\begin{theorem}[Joint convexity on positive-definite pairs]
    \label{thm:matrix_positive_definite_relative_entropy_convexity}
    \lean{TNLean.RelativeEntropyConvexity.convexOn_quantumRelativeEntropy_index}
    \leanok
    The function $(X,Y)\mapsto D_0(X\Vert Y)$ is jointly convex on pairs
    of positive-definite matrices. Equivalently,
    \begin{align}
        D_0(aX_1+bX_2\Vert aY_1+bY_2)
        \le aD_0(X_1\Vert Y_1)+bD_0(X_2\Vert Y_2)
        \label{eq:matrix_positive_definite_relative_entropy_convexity}
    \end{align}
    whenever $a,b\ge0$ and $a+b=1$. The affine trace correction in
    \cite[Fact~4]{Tropp2012JointConvexity} does not affect convexity.
\end{theorem}
\begin{proof}\leanok
    For $0\le s<1$, Lieb's trace-power concavity makes
    $(\operatorname{tr}X-\operatorname{tr}(X^sY^{1-s}))/(1-s)$ jointly convex.
    These functions converge to $D_0(X\Vert Y)$ as $s\uparrow1$.
    Convexity passes to the pointwise limit.
\end{proof}

\begin{theorem}[Lieb's exponential trace concavity]
    \label{thm:matrix_trace_exp_log_concavity}
    \lean{Matrix.concaveOn_trace_exp_add_log}
    \leanok
    \uses{thm:matrix_trace_variational_maximum,
        thm:matrix_positive_definite_relative_entropy_convexity}
    For every Hermitian $H$, the function
    \begin{align}
        A\longmapsto\operatorname{Re}\operatorname{tr}\exp(H+\log A)
        \label{eq:matrix_trace_exp_log_concavity}
    \end{align}
    is concave on the positive-definite cone. This is
    \cite[Theorem~1]{Tropp2012JointConvexity}.
\end{theorem}
\begin{proof}\leanok
    \uses{thm:matrix_trace_variational_maximum,
        thm:matrix_positive_definite_relative_entropy_convexity}
    Theorem~\ref{thm:matrix_trace_variational_maximum} writes the function as
    the maximum, over $X>0$, of
    $\operatorname{Re}\operatorname{tr}(XH)-D_0(X\Vert A)+
        \operatorname{Re}\operatorname{tr}X$.
    This expression is jointly concave, and its maximum at each $A$ is
    attained by $X=\exp(H+\log A)$. Mixing the two attaining matrices gives
    the asserted concavity.
\end{proof}

\begin{theorem}[Finite uniform trace-exponential Jensen inequality]
    \label{thm:finite_matrix_trace_exp_jensen}
    \lean{Matrix.average_trace_exp_add_le_trace_exp_log_average_exp}
    \leanok
    \uses{thm:matrix_trace_exp_log_concavity}
    Let $H$ be Hermitian and let $Y_a$, indexed by a finite nonempty set, be
    Hermitian. With uniform averaging denoted by $\mathbb E_a$,
    \begin{align}
        \mathbb E_a\operatorname{tr}e^{H+Y_a}
        \le\operatorname{tr}e^{H+\log(\mathbb E_a e^{Y_a})}.
        \label{eq:finite_matrix_trace_exp_jensen}
    \end{align}
    This is the finite uniform specialization of
    \cite[Corollary~3.3]{Tropp2012TailBounds}.
\end{theorem}
\begin{proof}\leanok
    \uses{thm:matrix_trace_exp_log_concavity}
    Apply Jensen's inequality on the positive-definite cone to $e^{Y_a}$.
\end{proof}

\begin{theorem}[Finite uniform logarithmic moment generating function bound]
    \label{thm:finite_matrix_log_mgf_bernstein}
    \lean{Matrix.log_average_exp_le_bernstein}
    \leanok
    Let $Z_a$ be a centered finite uniform Hermitian family with $\|Z_a\|\le R$.
    For $\theta\ge0$ and $\theta R<3$, put
    $V=\mathbb E_a Z_a^2$. Then
    \begin{align}
        \log\mathbb E_a e^{\theta Z_a}
        \le\frac{\theta^2\|V\|}{2(1-\theta R/3)}I.
        \label{eq:finite_matrix_log_mgf_bernstein}
    \end{align}
\end{theorem}
\begin{proof}\leanok
    The factorial inequality $2\,3^k\le(k+2)!$ bounds the exponential remainder
    by a geometric series. Functional calculus gives
    $e^{\theta Z_a}\le I+\theta Z_a+
        \theta^2Z_a^2/(2(1-\theta R/3))$.
    Averaging removes the linear term. The variance norm bounds the quadratic
    term by a scalar multiple of the identity, and $\log x\le x-1$ bounds
    each spectral value of the logarithm.
\end{proof}

\begin{theorem}[Finite uniform one-sided Hermitian Bernstein inequality]
    \label{thm:finite_uniform_hermitian_bernstein}
    \lean{Matrix.finite_uniform_hermitian_bernstein_of_variance_le}
    \leanok
    \uses{thm:finite_matrix_trace_exp_jensen,thm:finite_matrix_log_mgf_bernstein}
    Let $Z_a$ be a centered finite uniform family of Hermitian $d\times d$
    matrices with $\|Z_a\|\le R$. Let $\sigma^2\ge\|\mathbb E_a Z_a^2\|$.
    For $m$ independent uniform samples and $\tau\ge0$,
    \begin{align}
        \Pr\left(\lambda_{\max}\left(\sum_{j=1}^m Z_{a_j}\right)
        \ge\tau\right)
        \le d\exp\left(-\frac{\tau^2}{2m\sigma^2+2R\tau/3}\right).
        \label{eq:finite_uniform_hermitian_bernstein}
    \end{align}
    At a zero denominator the right side is read as $d$.
    This is the finite uniform, norm-bounded specialization of
    \cite[Theorem~6.1(ii)]{Tropp2012TailBounds}.
    The general probability theorem and the non-Hermitian norm estimate
    remain separate assertions.
\end{theorem}
\begin{proof}\leanok
    \uses{thm:finite_matrix_trace_exp_jensen,thm:finite_matrix_log_mgf_bernstein}
    Iterate the trace-exponential Jensen inequality over the sample coordinates.
    The logarithmic moment generating function bound and exponential Markov
    inequality give the usual scalar parameter bound. For $m\sigma^2>0$ take
    $\theta=\tau/(m\sigma^2+R\tau/3)$.
    If $m\sigma^2=0<R\tau$, take $\theta=3/(2R)$.
    The zero-denominator case follows by bounding a probability by one.
\end{proof}

\begin{lemma}[Hermitian dilation and variance domination]
    \label{lem:matrix_hermitian_dilation}
    \lean{Matrix.hermitianDilation,
        Matrix.isHermitian_hermitianDilation,
        Matrix.hermitianDilation_sq,
        Matrix.norm_hermitianDilation,
        Matrix.norm_mem_spectrum_hermitianDilation,
        Matrix.hermitianDilation_sum,
        Matrix.hermitianDilation_real_smul,
        Matrix.norm_average_hermitianDilation_sq_le}
    \leanok
    For a complex $d_1\times d_2$ matrix $Y$, define
    \begin{align}
        \mathcal D(Y)=\begin{pmatrix}0&Y\\Y^*&0\end{pmatrix}.
        \label{eq:matrix_hermitian_dilation}
    \end{align}
    This matrix is Hermitian, has norm $\|Y\|$, and has $\|Y\|$ in its spectrum.
    The dilation is real-linear and
    $\mathcal D(Y)^2=\operatorname{diag}(YY^*,Y^*Y)$.
    For a finite uniform family,
    \begin{align}
        \|\mathbb E\mathcal D(Y)^2\|
        \le\max\{\|\mathbb EYY^*\|,\|\mathbb EY^*Y\|\}.
        \label{eq:matrix_hermitian_dilation_variance}
    \end{align}
    These are the dilation properties used in
    \cite[equations~(2.12)--(2.13), Remark~3.11]{Tropp2012TailBounds}.
\end{lemma}
\begin{proof}\leanok
    The square identity and positive-semidefinite block domination bound the
    dilation norm from above. The two coordinate inclusions bound it from below.
    Conjugation by $\operatorname{diag}(I,-I)$ changes the sign of the dilation,
    so its spectrum is symmetric and contains its positive norm endpoint.
    The same block domination gives the variance bound.
\end{proof}

\begin{theorem}[Finite uniform rectangular matrix Bernstein inequality]
    \label{thm:finite_uniform_rectangular_bernstein}
    \lean{Matrix.finite_uniform_rectangular_bernstein,
        Matrix.finite_uniform_rectangular_average_bernstein}
    \leanok
    \uses{lem:matrix_hermitian_dilation,thm:finite_uniform_hermitian_bernstein}
    Let $Y$ be a centered finite uniform family of complex $d_1\times d_2$
    matrices, with $\|Y\|\le R$, and put
    $\sigma^2=\max\{\|\mathbb EYY^*\|,\|\mathbb EY^*Y\|\}$.
    For $\ell\ge1$ independent uniform samples and $t\ge0$,
    \begin{align}
        \Pr\left(\left\|\frac1\ell\sum_{j=1}^{\ell}Y_j\right\|\ge t\right)
        \le(d_1+d_2)\exp\left(-\frac{\ell t^2}{2\sigma^2+2Rt/3}\right).
        \label{eq:finite_uniform_rectangular_bernstein}
    \end{align}
    This is the finite uniform specialization of
    \cite[Theorem~1.6]{Tropp2012TailBounds}.
\end{theorem}
\begin{proof}\leanok
    \uses{lem:matrix_hermitian_dilation,thm:finite_uniform_hermitian_bernstein}
    Apply the one-sided Hermitian theorem to the dilation and use its positive
    spectral norm endpoint. The dilation dimension is $d_1+d_2$.
    Normalize the sum by $\ell$.
\end{proof}

\begin{corollary}[Finite uniform concentration about the actual matrix mean]
    \label{cor:finite_uniform_square_matrix_bernstein}
    \lean{Matrix.finite_uniform_matrix_average_sub_mean_bernstein}
    \leanok
    \uses{thm:finite_uniform_rectangular_bernstein}
    Let $X$ be a finite uniform family of complex $d\times d$ matrices and
    $\mu=\mathbb EX$. If $\|X-\mu\|\le R$, set
    $\sigma^2=\max\{\|\mathbb E(X-\mu)(X-\mu)^*\|,
        \|\mathbb E(X-\mu)^*(X-\mu)\|\}$.
    The bound in Theorem~\ref{thm:finite_uniform_rectangular_bernstein}
    holds for the deviation of the sample mean from $\mu$, with factor $2d$.
    This is the finite uniform specialization of LVV Theorem~2.13.
    The version for general probability spaces is
    Theorem~\ref{thm:lvv_matrix_bernstein}.
\end{corollary}
\begin{proof}\leanok
    \uses{thm:finite_uniform_rectangular_bernstein}
    Subtract the actual uniform mean. The centered sum and its normalization
    are computed explicitly, then the rectangular theorem is applied with
    equal row and column dimensions.
\end{proof}

\begin{lemma}[Integrability of continuous functions of bounded matrix families]
    \label{lem:integrable_continuous_bounded_matrix_family}
    \lean{Matrix.integrable_continuous_comp_of_ae_norm_le}
    \leanok
    On a finite measure space, let $Y$ be an almost everywhere bounded strongly
    measurable finite matrix family. For every continuous function $g$ into a
    normed additive group, $g\circ Y$ is integrable.
\end{lemma}
\begin{proof}\leanok
    A closed matrix ball is compact. Its continuous image is bounded, which
    gives an integrable constant bound for $\|g\circ Y\|$.
\end{proof}

\begin{theorem}[Trace-exponential Jensen inequality for bounded random matrices]
    \label{thm:bounded_matrix_trace_exp_integral_jensen}
    \lean{Matrix.integral_trace_exp_add_le_trace_exp_log_integral_exp}
    \leanok
    \uses{thm:matrix_trace_exp_log_concavity,
        lem:integrable_continuous_bounded_matrix_family}
    Let $H$ be Hermitian. On a probability space, let $Y$ be strongly measurable,
    Hermitian almost everywhere, and bounded in operator norm almost everywhere.
    Then
    \begin{align}
        \mathbb E\operatorname{tr}e^{H+Y}
        \le\operatorname{tr}e^{H+\log\mathbb Ee^Y}.
        \label{eq:bounded_matrix_trace_exp_integral_jensen}
    \end{align}
    This is the bounded-sample specialization of
    \cite[Corollary~3.3]{Tropp2012TailBounds}.
\end{theorem}
\begin{proof}\leanok
    \uses{thm:matrix_trace_exp_log_concavity,
        lem:integrable_continuous_bounded_matrix_family}
    If $\|Y\|\le R$, functional calculus gives $e^Y\ge e^{-R}I$.
    Restrict Lieb's concavity to the closed convex set of matrices at least
    $e^{-R}I$. The logarithm is continuous there. The preceding integrability
    lemma verifies the hypotheses of integral Jensen; $\log e^Y=Y$ gives
    the stated inequality.
\end{proof}

\begin{theorem}[Integral logarithmic moment generating function bound]
    \label{thm:bounded_matrix_integral_log_mgf_bernstein}
    \lean{Matrix.log_integral_exp_le_bernstein}
    \leanok
    Let $Z$ be a strongly measurable, centered Hermitian random matrix on a
    probability space, with $\|Z\|\le R$ almost everywhere. For $\theta\ge0$
    and $\theta R<3$,
    \begin{align}
        \log\mathbb Ee^{\theta Z}
        \le\frac{\theta^2\|\mathbb EZ^2\|}{2(1-\theta R/3)}I.
        \label{eq:bounded_matrix_integral_log_mgf_bernstein}
    \end{align}
\end{theorem}
\begin{proof}\leanok
    \uses{lem:integrable_continuous_bounded_matrix_family}
    The scalar Bernstein majorant and functional calculus give the quadratic
    exponential bound. Integrate in the Loewner order and use centering.
    Bound the variance by its norm times the identity, then apply operator
    monotonicity of the logarithm and $\log(1+x)\le x$.
    The positive lower bound $e^{\theta Z}\ge e^{-\theta R}I$ ensures that
    the expected exponential is positive definite.
\end{proof}

\begin{lemma}[Almost-everywhere bounds and integrability for finite matrix sums]
    \label{lem:bounded_matrix_sum_integrability}
    \lean{Matrix.ae_norm_sum_le_of_forall,
        Matrix.integrable_trace_exp_sum_of_ae_norm_le}
    \leanok
    Let $X_1,\ldots,X_m$ be almost everywhere measurable matrix families on a
    finite measure space, with $\|X_j\|\le R$ almost everywhere. Then
    $\|\sum_jX_j\|\le mR$ almost everywhere. For every fixed matrix $H$,
    $\operatorname{Re}\operatorname{tr}e^{H+\sum_jX_j}$ is integrable.
    Independence and Hermiticity are unnecessary for this conclusion.
\end{lemma}
\begin{proof}\leanok
    \uses{lem:integrable_continuous_bounded_matrix_family}
    Intersect the finitely many events of full measure and use the triangle
    inequality. Apply the continuous-function integrability lemma to the sum.
\end{proof}

\begin{theorem}[Trace-exponential iteration on a product probability space]
    \label{thm:bounded_product_matrix_trace_exp}
    \lean{Matrix.integral_trace_exp_sum_pi_le}
    \leanok
    \uses{thm:bounded_matrix_trace_exp_integral_jensen,
        lem:bounded_matrix_sum_integrability}
    Let $\mu_1,\ldots,\mu_m$ be probability distributions on Hermitian matrices,
    each bounded in norm by $R$ almost everywhere. On their product space,
    for every Hermitian $H$,
    \begin{align}
        \mathbb E\operatorname{tr}e^{H+\sum_jX_j}
        \le\operatorname{tr}e^{H+\sum_j\log\mathbb E_{\mu_j}e^{X_j}}.
        \label{eq:bounded_product_matrix_trace_exp}
    \end{align}
    The marginal distributions may differ.
\end{theorem}
\begin{proof}\leanok
    \uses{thm:bounded_matrix_trace_exp_integral_jensen,
        lem:bounded_matrix_sum_integrability}
    Separate the first coordinate by the measurable product equivalence and
    apply Fubini. Integral Jensen replaces that coordinate by its logarithmic
    moment generating function. Induct on the number of coordinates, allowing
    the fixed Hermitian shift to change at each step.
\end{proof}

\begin{theorem}[Trace-exponential bound for bounded independent random matrices]
    \label{thm:bounded_independent_matrix_trace_exp}
    \lean{Matrix.integral_trace_exp_sum_le_of_iIndepFun}
    \leanok
    \uses{thm:bounded_product_matrix_trace_exp}
    Let $X_1,\ldots,X_m$ be independent, almost everywhere measurable Hermitian
    random matrices on a probability space, each bounded in norm by $R$ almost
    everywhere. For every Hermitian $H$,
    \begin{align}
        \mathbb E\operatorname{tr}e^{H+\sum_jX_j}
        \le\operatorname{tr}e^{H+\sum_j\log\mathbb Ee^{X_j}}.
        \label{eq:bounded_independent_matrix_trace_exp}
    \end{align}
    This is the bounded form of \cite[Lemma~3.4]{Tropp2012TailBounds} used in
    the Bernstein argument.
\end{theorem}
\begin{proof}\leanok
    \uses{thm:bounded_product_matrix_trace_exp}
    Independence identifies the actual joint distribution with the product
    of the marginal distributions. Change variables in the joint integral
    and in each marginal exponential integral, then apply the product-space
    theorem.
\end{proof}

\begin{theorem}[Trace-exponential estimate from an upper variance bound]
    \label{thm:independent_matrix_trace_exp_variance_bound}
    \lean{Matrix.integral_trace_exp_sum_le_bernstein_of_variance_le}
    \leanok
    \uses{thm:bounded_independent_matrix_trace_exp,
        thm:bounded_matrix_integral_log_mgf_bernstein}
    Let $X_1,\ldots,X_m$ be independent centered Hermitian random matrices on a
    probability space, with $\|X_j\|\le R$ almost everywhere and
    $\|\mathbb EX_j^2\|\le\sigma^2$ for every $j$.
    For $\theta\ge0$ and $\theta R<3$,
    \begin{align}
        \mathbb E\operatorname{tr}e^{\theta\sum_jX_j}
        \le d\exp\left(\frac{m\theta^2\sigma^2}{2(1-\theta R/3)}\right).
        \label{eq:independent_matrix_trace_exp_variance_bound}
    \end{align}
\end{theorem}
\begin{proof}\leanok
    \uses{thm:bounded_independent_matrix_trace_exp,
        thm:bounded_matrix_integral_log_mgf_bernstein}
    The independent trace-exponential inequality bounds the left side by the
    trace exponential of the sum of the logarithmic moment generating
    functions. Bound each summand by
    $\theta^2\sigma^2I/(2(1-\theta R/3))$ and use monotonicity of the trace
    exponential.
\end{proof}

\begin{theorem}[Exponential Markov bound for the upper spectral event]
    \label{thm:independent_matrix_spectral_parameter_bound}
    \lean{Matrix.probability_spectral_tail_le_bernstein_parameter}
    \leanok
    \uses{thm:independent_matrix_trace_exp_variance_bound}
    Under the preceding hypotheses, for every real $\tau$,
    \begin{align}
        \Pr\left(\lambda_{\max}\left(\sum_jX_j\right)\ge\tau\right)
        \le d\exp\left(-\theta\tau+
        \frac{m\theta^2\sigma^2}{2(1-\theta R/3)}\right).
        \label{eq:independent_matrix_spectral_parameter_bound}
    \end{align}
\end{theorem}
\begin{proof}\leanok
    \uses{thm:independent_matrix_trace_exp_variance_bound}
    The spectral event implies
    $e^{\theta\tau}\le\operatorname{tr}e^{\theta\sum_jX_j}$ almost surely.
    Apply scalar exponential Markov to this actual trace-exponential random
    variable. Its integrability and nonnegativity have been established.
\end{proof}

\begin{theorem}[One-sided Hermitian Bernstein with an upper variance bound]
    \label{thm:independent_hermitian_bernstein_variance_bound}
    \lean{Matrix.probability_hermitian_bernstein_of_variance_le}
    \leanok
    \uses{thm:independent_matrix_spectral_parameter_bound}
    Let $X_1,\ldots,X_m$ be independent centered Hermitian random matrices on a
    probability space. Suppose $\|X_j\|\le R$ almost everywhere and
    $\|\mathbb EX_j^2\|\le\sigma^2$, where $R,\sigma^2\ge0$.
    For $\tau\ge0$,
    \begin{align}
        \Pr\left(\lambda_{\max}\left(\sum_jX_j\right)\ge\tau\right)
        \le d\exp\left(-\frac{\tau^2}{2m\sigma^2+2R\tau/3}\right).
        \label{eq:independent_hermitian_bernstein_variance_bound}
    \end{align}
    This is the upper-variance-bound consequence of the argument in
    \cite[Theorem~6.1(ii)]{Tropp2012TailBounds}. It uses $m\sigma^2$ in place
    of the sharper norm of the sum of variances for unrelated distributions.
\end{theorem}
\begin{proof}\leanok
    \uses{thm:independent_matrix_spectral_parameter_bound}
    If $m\sigma^2>0$, take
    $\theta=\tau/(m\sigma^2+R\tau/3)$.
    If $m\sigma^2=0<R\tau$, take $\theta=3/(2R)$.
    In the remaining zero-denominator case, take $\theta=0$ and use the
    inverse-zero convention. Thus zero variance, threshold and sample count
    are included without an additional hypothesis.
\end{proof}

\begin{lemma}[Hermitian dilation and the Bochner integral]
    \label{lem:matrix_integral_hermitian_dilation}
    \lean{Matrix.hermitianDilationLinearIsometry,
        Matrix.integral_hermitianDilation,
        Matrix.norm_integral_hermitianDilation_sq_le}
    \leanok
    \uses{lem:matrix_hermitian_dilation}
    The Hermitian dilation of a complex square matrix is a real-linear isometry.
    It therefore commutes with the Bochner integral. For an almost everywhere
    measurable bounded matrix family $Y$ on a finite measure space,
    \begin{align}
        \left\|\int\mathcal D(Y)^2\right\|
        \le\max\left\{\left\|\int YY^*\right\|,
        \left\|\int Y^*Y\right\|\right\}.
        \label{eq:matrix_integral_hermitian_dilation_variance}
    \end{align}
\end{lemma}
\begin{proof}\leanok
    \uses{lem:matrix_hermitian_dilation}
    Real-linearity and preservation of the operator norm give the isometry.
    Its integral identity follows from linear isometry of the Bochner integral.
    The two continuous linear corner embeddings identify the integrated squared
    dilation with the block diagonal of the integrated Gram matrices.
    These blocks are positive semidefinite; domination by the maximum of their
    norms times the identity gives the variance bound.
\end{proof}

\begin{theorem}[The iid complex matrix Bernstein inequality]
    \label{thm:lvv_matrix_bernstein}
    \lean{Matrix.probability_matrix_bernstein}
    \leanok
    \uses{thm:independent_hermitian_bernstein_variance_bound,
        lem:matrix_integral_hermitian_dilation}
    Let $X_1,\ldots,X_\ell$ be independent identically distributed complex
    $d\times d$ random matrices on a probability space, with $d,\ell\ge1$,
    common mean $\mu$ and $\|X_j-\mu\|\le R$ almost surely. Put $Y=X_1-\mu$ and
    \begin{align}
        \sigma^2=\max\{\|\mathbb EYY^*\|,\|\mathbb EY^*Y\|\}.
        \label{eq:lvv_matrix_bernstein_variance}
    \end{align}
    For every $t\ge0$,
    \begin{align}
        \Pr\left(\left\|\frac1\ell\sum_{j=1}^{\ell}X_j-\mu\right\|\ge t\right)
        \le2d\exp\left(-\frac{\ell t^2}{2\sigma^2+2Rt/3}\right).
        \label{eq:lvv_matrix_bernstein}
    \end{align}
    This is LVV Theorem~2.13, equation~(7), the iid complex square-matrix case
    of \cite[Theorem~1.6]{Tropp2012TailBounds}. The scalar maximum of variance
    norms clarifies the notation in LVV. At a zero denominator the inverse-zero
    convention gives the bound $2d$.
\end{theorem}
\begin{proof}\leanok
    \uses{thm:independent_hermitian_bernstein_variance_bound,
        lem:matrix_integral_hermitian_dilation}
    Derive integrability from the deviation bound and center the actual samples.
    Apply the Hermitian dilation to each centered matrix. Its mean is zero and
    its norm has the same bound. Identical distribution identifies both Gram
    expectations with those of the first actual sample; the preceding lemma
    supplies the variance cap. Apply the one-sided Hermitian bound at threshold
    $\ell t$. The positive spectral norm endpoint gives dimension $2d$.
    The explicit centering and sample-mean identities give the stated event
    and exponent.
\end{proof}

Theorem~\ref{thm:lvv_matrix_bernstein} supplies the sharp concentration result.
The sampled-filter construction in LVV Corollary~2.14 also requires a
simultaneous word-frequency bound and the corrected ground-energy dependence;
these are separate assertions.
